\section{A primitive recursive Mordell--Weil descent schema}
\label{sec:descent}

Let $\Gamma$ be an abelian group represented by valid natural-number codes,
with fixed primitive recursive equality, addition, negation and scalar
multiplication. Suppose its laws are proved on valid codes in iPRA.
The group is not required to be finite or to have a decidable rank.
Let $h:\Gamma\to\N$ be a fixed primitive recursive height.

An enumeration $L(B)$ means a primitive recursive finite list of valid
points, with the proved property
\[
 h(P)\le B\quad\Longrightarrow\quad P\text{ occurs in }L(B).
 \tag{3.1}\label{eq:enum}
\]
Duplicates are harmless. Including the identity in every list avoids
irrelevant empty-list branches.

\begin{theorem}[Explicit contracting descent]\label{thm:abstract}
Suppose $R=(R_0,\ldots,R_{r-1})$, $r>0$, is a finite list, $C\in\N$,
and fixed primitive recursive functions $i(P),D(P)$ satisfy, for all
valid $P$ with $h(P)>C$,
\[
 i(P)<r,\qquad D(P)\in\Gamma,\qquad
 P=2D(P)+R_{i(P)},\qquad
 2h(D(P))\le h(P)+C.
 \tag{3.2}\label{eq:contract}
\]
Assume \eqref{eq:enum} and these statements have iPRA proofs.
Then $G=R\mathbin{+\!\!+}L(C)$ generates $\Gamma$. A primitive recursive
function produces the coefficients in \eqref{eq:1-coord} below, with an
iPRA proof of correctness.

More explicitly, write $e=(h(P)\dotminus C)$ and let
$t=\bl(e)$, where $\bl(0)=0$ and, for $e>0$, $\bl(e)$ is the unique
integer with $2^{t-1}\le e<2^t$. The algorithm takes at most $t$ active
descent steps. Every resulting coefficient is at most $2^t$; if $e>0$,
$2^t\le2e$, and if $e=0$ the largest coefficient is $1$.
\end{theorem}

\begin{proof}
Define the excess $e(Q)=h(Q)\dotminus C$. An active step has
$h(Q)=C+e(Q)$, so \eqref{eq:contract} gives
$2h(D(Q))\le2C+e(Q)$. Therefore
\[
 e(D(Q))\le \left\lfloor e(Q)/2\right\rfloor.
 \tag{3.3}\label{eq:halve-excess}
\]
If the child has height below $C$, the left side is zero; otherwise this
is direct subtraction. In particular every active step strictly decreases
height. Freeze the whole state as soon as its height is at most $C$.

The state consists of a current point $Q$, an active-step count $k$,
and an $r$-entry coefficient vector $a$. Start at
$Q=P,k=0,a=0$. An active step, using the old point, performs
\[
 j=i(Q),\quad Q'=D(Q),\quad
 a'_j=a_j+2^k,\quad a'_l=a_l\ (l\ne j),\quad k'=k+1.
 \tag{3.4}\label{eq:3-update}
\]
Run this stopped transition for exactly $t$ stages. The natural number $t$
is computed by a bounded scan: search $0,\ldots,e+1$ for the least $j$
with $e<2^j$. Induction proves $2^{e+1}>e$, so the scan succeeds.

After $j$ stages, \eqref{eq:halve-excess} gives
$e(Q)\le\lfloor e/2^j\rfloor$. A convenient entirely integral invariant
is $2^j e(Q)\le e$, maintained also through frozen stages because then
$e(Q)=0$. At $j=t$, $e<2^t$, so $e(Q)=0$ and $Q$ belongs to $L(C)$.
If one prefers a simpler bound, $h(P)$ stages suffice by strict decrease.
The compiled Lean kernel uses this coarser bound.

The simultaneous group invariant is
\[
 P=2^kQ+\sum_{l<r}a_lR_l.
 \tag{3.5}\label{eq:group-invariant}
\]
It holds initially. Substituting $Q=2Q'+R_j$ proves it after
\eqref{eq:3-update}; a frozen step leaves it unchanged.
At the end, scan $L(C)$ for $Q$ using decidable equality and
Lemma~\ref{lem:scan}. If $Q=L(C)_b$, then
\[
 P=\sum_{l<r}a_lR_l+2^kL(C)_b.
 \tag{3.6}\label{eq:1-coord}
\]
The digit weights added in active steps are $1,2,\ldots,2^{k-1}$.
Consequently $\sum_l a_l=2^k-1$, and $k\le t$ proves the coefficient
bound. The elementary inequality $2^t\le2e$ for $e>0$ follows from the
minimality of $t$.

For the iPRA proof, encode the vector and state. Validity, the excess
inequality, the coefficient-sum identity and
\eqref{eq:group-invariant} are quantifier-free formulas about the coded
functions. The finite sum is a primitive recursive fold. Correctness of
each update supplies the induction step of \eqref{eq:qfind}. The possibly
quantified global premises remain hypotheses of the derivation; they
are not placed inside its induction formula.
\end{proof}

\subsection{A formulation requiring no group identities}

The core also works for a set with a binary operation $\oplus$, without
associativity or commutativity. Assume a child satisfies
$P=(D(P)\oplus D(P))\oplus R_{i(P)}$.
Construct a word by returning the corresponding bounded-height atom when
small, and otherwise returning
\[
 \operatorname{add}\bigl(\operatorname{double}(w(D(P))),
                         \operatorname{rep}(i(P))\bigr).
 \tag{3.7}
\]
With $h(P)$ fuel this gives an expression evaluating exactly to $P$.
Each active stage adds three constructors, so the compact expression has
at most $3h(P)+1$ constructors. A dedicated \emph{double} constructor
avoids physically copying the child expression.

This is the theorem proved in \code{lean/DescentCore.lean}. In an abelian
group, folding the word into coefficients gives (3.6). For elliptic curves,
the word version isolates the arithmetic chord-and-tangent construction
from the independent proof of all abstract group identities.

\subsection{What counts as a certificate}

A finite list, a constant and program definitions do not make their
universal correctness automatically decidable. The assumptions in
Theorem~\ref{thm:abstract} are mathematical assertions with proofs.
There are two honest uses of the theorem: prove these assertions for a
particular input, or state a conditional theorem with those assertions as
premises.

One may also supply finite formal derivations of the premises. Checking
such a derivation is a syntactic operation, and combining it with a fixed
proof template can produce a new derivation. This does not justify an
internal principle saying that every accepted proof has a true conclusion.
Such a principle would be a reflection assertion and is not used here.

\subsection{A simple infinite test group}

For $\Gamma=\Z^r$, take $h(z)=\sum_j z_j^2$, digits
$(z_j\bmod2)_j\in\{0,1\}^r$, and child
$q_j=\lfloor z_j/2\rfloor$. For an even coordinate,
$2q_j^2\le z_j^2+1$; for $z_j=2q_j+1$ the same inequality follows from
$(2q_j+1)^2+1-2q_j^2=2(q_j+1)^2\ge0$. Summing gives the contract
with $C=r$. The finite cube $[-r,r]^r$, filtered by $\sum_jz_j^2\le r$,
enumerates the base region. This illustrates an infinite group where the
algorithm's hypotheses are elementary and the output is not vacuous.
It is an algorithmic test case, not an elliptic-curve rank computation.
